\section{Recovering the usual local hypothesis}\label{sec:local}

\subsection{Choosing the neighborhood from a base-point certificate}

Suppose the original \(F\) is quantitatively \(C^1\) on a
neighborhood containing a closed product box of radius \(R_0>0\)
about a zero \((a,b)\). Suppose rational matrices \(A,B\) satisfy
\[
 AB=BA=\Id,\qquad
 \norm{\Id-A D_yF(a,b)}\le\tfrac14.
 \tag{Base inverse certificate}
\]
Set \(H=AF\). Multiplying a value or a derivative evaluator by
\(A\), with an additional precision shift
\(\lgup(\max(1,\norm A))\), produces evaluators for \(H\) and
\(D H\). A Taylor or derivative continuity modulus for \(F\)
transfers by the same shift in its demanded error index.
We may therefore work with the resulting moduli for \(H\).

Put
\[
 R=\min(R_0/2,\,2^{-\nu(2)}),\qquad r=R.
\]
Throughout the product box of radius \(R\),
\(\norm{V(z)-V(a,b)}\le1/4\), so
\(\norm{\Id-V(z)}\le1/2\).
Evaluate \(U(a,b)\) to error \(1/4\), and choose the rational
bound
\[
 M=\max\bigl(1,\norm{\widehat U((a,b),2)}+\tfrac12\bigr).
\]
It bounds \(\norm U\) on this box: the evaluation error and
the variation from the base point each cost \(1/4\).
Now choose
\[
 \alpha=\frac{r}{4M}.
\]
Then \(\alpha\le r\), all points of \(X\times Y\) lie in the
radius-\(R\) box, and \(M\alpha=r/4\).
The exact equation \(H(a,b)=0\) satisfies the remaining
hypothesis. Theorem~\ref{thm:main} and
Corollary~\ref{cor:original} yield a unique local \(C^1\) graph.
All choices are rational operations or evaluations of supplied
primitive recursive moduli.

\subsection{Constructing the rational preconditioner}
\label{sec:conditioning}

Here is a precise sufficient replacement for an unspecified
invertibility assertion. Let \(J_0=D_yF(a,b)\). Suppose
\[
 |(J_0)_{ij}|\le K,\quad K\ge1,\qquad
 |\det J_0|\ge\delta>0,
\]
where \(K,\delta\) are rational and the bounds are certified.
Define
\[
 D=m!(K+1)^{m-1},\qquad
 \eta=\min\left(1,\frac{\delta}{2mD},
                         \frac{\delta}{8D}\right)>0.
\]
Choose a rational matrix \(R\) approximating \(J_0\) in row norm
to error at most \(\eta\), by evaluating at precision
\(\lgup(1/\eta)\).

For each determinant product, telescope its \(m\) factors.
Every factor has absolute value at most \(K+1\), and each
factor changes by at most \(\eta\). Summing over \(m!\)
permutations gives the finite estimate
\[
 |\det R-\det J_0|
 \le m!m(K+1)^{m-1}\eta=mD\eta\le\delta/2.
\]
It follows that the rational number \(\det R\) satisfies
\(|\det R|\ge\delta/2>0\). Its nonzero denominator and nonzero
determinant can be checked by rational arithmetic.
Define
\[
 A=\frac{\operatorname{adj}R}{\det R},\qquad B=R.
\]
The finite adjugate identities prove \(AB=BA=\Id\).
Each cofactor is bounded by \((m-1)!(K+1)^{m-1}\); summing a row
of \(m\) cofactors gives
\[
 \norm A\le\frac{2D}{\delta}.
\]
Therefore
\[
 \norm{\Id-AJ_0}
 =\norm{A(R-J_0)}
 \le\frac{2D}{\delta}\eta\le\tfrac14.
\]
This is the base inverse certificate used above. No search for
a successful matrix inverse occurs: the evaluation accuracy is
chosen before the matrix is computed.

The argument uses determinant \emph{apartness}, a positive
lower-bound witness. If a constructive inverse name and a bound
on its norm are given instead, approximating that inverse also
produces suitable quantitative data. A bare negative proposition
about a determinant is not used to manufacture a numerical
witness.

\subsection{Base points given by names}

Rational centers are a convenience for the rounded program,
not a restriction to rational zeros. If \(a,b\) have fixed
primitive recursive names, work in translated coordinates:
\[
 F^\circ(u,v)=F(a+u,b+v).
\]
The new box has rational center \((0,0)\).
Its evaluator is obtained by approximating \(a,b\) with a
precision prescribed by the local continuity bound for \(F\),
then evaluating \(F\) at that rational approximation.
For \(D F^\circ\), use the supplied derivative modulus.
Choose the translated box strictly inside a larger certified
domain so that the approximating rational points remain in
that domain. The \(R_0/2\) margin above supplies this slack.

For example, if \(F\) has Lipschitz bound \(K_F\) on the larger
box, use the translated-input precision
\[
 t=\max\bigl(p+1+\lgup(\max(1,K_F)),\ \lgup(2/R_0)\bigr).
\]
The second term ensures that the input error is at most
\(R_0/2\), so the rational evaluation stays in the larger
domain. Evaluating \(F\) to \(2^{-(p+1)}\) then gives total
value error at most \(2^{-p}\).
Derivative evaluation uses its modulus-driven precision
with the same additional domain guard. Translation preserves the Taylor inequalities,
and uniform slack enlargements of the moduli handle rational
approximations exactly as in Section~\ref{sec:raw}.
Apply the theorem to \(F^\circ\), obtaining \(v=g^\circ(u)\);
the original branch is \(b+g^\circ(x-a)\).

\subsection{The constructive inverse function theorem}

Let \(\Phi\) be a quantitatively \(C^1\) map between equal
finite dimensions, with positively conditioned derivative at
\(b\). Set \(a=\Phi(b)\) and
\[
 F(x,y)=\Phi(y)-x.
\]
Then \(D_yF=D\Phi\), and after preconditioning
\(D_xH=-A\), so one can take \(M\ge\norm A\).
The preceding construction gives a primitive recursive local
inverse \(g\) with
\[
 \Phi(g(x))=x,\qquad
 g(\Phi(y))=y
\]
whenever the displayed points remain in the constructed
neighborhoods. The second identity follows from uniqueness.
Moreover
\[
 Dg(x)=D\Phi(g(x))^{-1},
\]
with the derivative algorithm and explicit moduli already given.
Thus the finite construction also supplies the usual local
coordinate-change consequence.

\subsection{Scope of the theorem}

The proved result is finite-dimensional \(C^1\) regularity.
It does not claim the full Banach-space theorem, an analytic
implicit-function theorem, or automatic modulus extraction
from arbitrary qualitative input. Fixed higher differentiability
orders can be approached by differentiating the displayed
matrix formula with additional derivative certificates, but
those extensions are not needed for, or included in, the
proved theorem.
